\chapter{Составная проекторная связность и граница спинорного накрытия}

Глобальный проекторный ёж имеет правильный класс $+15/-15$, но его
квадратичная энергия расходится. Стандартный ремонт вводит независимую
$SO(3)$-связность. Однако прежде чем расширять теорию новым калибровочным полем,
необходимо проверить, не содержится ли требуемая связность в самом
пространственно меняющемся параметре порядка $Q$.

\section{Каноническая составная связность}

Пусть вакуумная щель упорядоченной фазы равна
\begin{equation}
 \Delta=a-b=0.8682499005\ldots.
\end{equation}
Определим антисимметричную матричную одноформу
\begin{equation}
 A_Q:=\frac{[Q,dQ]}{\Delta^2}.
 \label{eq:v6-composite-projector-connection}
\end{equation}
Для одноосного поля
\begin{equation}
 Q=q\left(P-\frac{I_3}{3}\right),
 \qquad P=nn^T,
\end{equation}
получаем
\begin{equation}
 A_Q=\left(\frac q\Delta\right)^2[P,dP].
 \label{eq:v6-composite-uniaxial-connection}
\end{equation}
Новая непрерывная константа здесь не введена: знаменатель фиксирован уже
вычисленной вакуумной щелью.

Тождество для проектора
\begin{equation}
 [[P,dP],P]=-dP
\end{equation}
показывает, что при $q\to\Delta$ ковариантная производная
\begin{equation}
 D_QQ=dQ+[A_Q,Q]
\end{equation}
точно удаляет дальнее вращение оси. Тем самым линейная инфракрасная
расходимость глобального ежа исчезает без независимо вставленной
пространственной связности.

Формулы такого типа родственны restricted connection Чо и переменным
Фаддеева--Ниеми; см. DOI \path{10.1103/PhysRevD.21.1080}, DOI
\path{10.1103/PhysRevLett.82.1624} и
\path{arXiv:hep-th/9807069}. В проекте конструкция возникает не как
перепараметризация готового Yang--Mills поля, а непосредственно из $Q$.

\section{Гладкий конечный пробный профиль}

Для ежа выберем
\begin{equation}
 P(x)=\frac{xx^T}{r^2},
 \qquad
 q(r)=\Delta\frac{r^2}{r^2+R^2}.
 \label{eq:v6-composite-smooth-profile}
\end{equation}
Тогда $Q$ гладок в начале координат, поскольку $qP$ пропорционален
$xx^T$ при $r\to0$. Связность имеет поведение
\begin{equation}
 A_Q=O(r^3),
\end{equation}
и также гладка на ядре.

Численный аудит при $R=1$ и единичных коэффициентах получил
\begin{align}
 \int|D_QQ|^2&=12.4786366\ldots,\\
 \int|F_{A_Q}|^2&=33.5371715\ldots,\\
 \int(\Delta^2-q^2)^2&=17.5002221\ldots.
 \label{eq:v6-composite-finite-integrals}
\end{align}
Дальние плотности убывают соответственно как
$r^{-6}$, $r^{-4}$ и $r^{-4}$. Все три радиальных интеграла сходятся.

\begin{theorem}[Составное устранение инфракрасной расходимости]
Поле $Q$ канонически порождает связность $A_Q=[Q,dQ]/\Delta^2$, для
которой существует гладкий проекторный ёж с конечными интегралами
$|D_QQ|^2$, $|F_{A_Q}|^2$ и объёмного потенциала.
\end{theorem}

Это сильнее предыдущего условного BPS-примера: существование конечной
конфигурации теперь использует только уже выведенное поле $Q$, его
вакуумную щель и проекторную кривизну.

\section{Почему это ещё не новое калибровочное поле}

Проверим локальный закон преобразования. Для
$Q\mapsto gQg^{-1}$ неоднородная часть преобразования
\eqref{eq:v6-composite-projector-connection} равна
\begin{equation}
 \frac{[Q,[d\omega,Q]]}{\Delta^2}.
 \label{eq:v6-composite-local-transform}
\end{equation}
На одноосном вакууме оператор $\operatorname{ad}_Q^2/\Delta^2$ является
проектором на два направления, меняющие ось. Поэтому для них возникает
правильный член $-d\omega$. Но генератор стабилизатора $O(2)$ коммутирует
с $Q$, и его неоднородная часть равна нулю.

Машинный тест дал остаток $1.6\cdot10^{-16}$ на двух сломанных
генераторах и полный пропуск стабилизаторного генератора.

\begin{proposition}[Связность косета, а не полный калибровочный бозон]
$A_Q$ является канонической составной связностью направлений
$SO(3)/O(2)$. Она глобально $SO(3)$-ковариантна и компенсирует вращение
директора, но из одного $Q$ не следует полный локальный закон
$SO(3)$-связности.
\end{proposition}

Следовательно, две моды Голдстоуна не превратились автоматически в три
независимых калибровочных бозонов. Мы получили геометрическую обратную связь поля
порядка, а не новый фундаментальный калибровочный сектор.

\section{Роль спинорного накрытия}

Недостающая стабилизаторная часть имеет правильный тип: на сфере вокруг
спинорного накрытия дефекта из Тома V уже даёт хопфову пару
\begin{equation}
 L\oplus L^*,
 \qquad c_1(L)=1,
 \qquad c_1(L^*)=-1,
\end{equation}
а коэффициентная кратность переводит её в классы $+15/-15$.
Таким образом, составная косетная связность и граничная хопфова линия
дополняют друг друга по типам.

Но прежний ранговый запрет не исчезает. Граница знает $L/L^*$, тогда как
гладкое ядро требует квадратного комплексно-линейного оператора,
обратимого вне центра. Прямой угол $20\times15$ оставляет постоянное
коядро размерности пять, а канонический равногранный ремонт Томом V не
найден.

Поэтому индекс Каллиаса ещё нельзя назначить новому конечному профилю.
Следующая проверка должна проверить, делает ли совместная пара
$(A_Q,L\oplus L^*)$ соответствующий оператор Фредгольмовым, или старый
ранговый запрет окончательно запрещает единичную локализованную моду.

\section{Стабильность размера}

При масштабировании $R$ три положительных вклада ведут себя как
\begin{equation}
 E_D\sim c_1R,
 \qquad
 E_F\sim\frac{c_2}{R},
 \qquad
 E_V\sim c_3R^3.
\end{equation}
Поэтому положительный кривизный член предотвращает коллапс и вместе с
объёмным потенциалом допускает конечный минимум. Но общий относительный
вес этих трёх членов всё ещё не выведен одним родительским следом; масса и
радиус остаются условными.

\begin{nogocriterion}[Составная связность не равна фундаментальной]
Конечность пробной энергии не позволяет объявить $A_Q$ новым калибровочным полем,
пока не доказаны полный локальный закон преобразования, кинетическая
нормировка и Фредгольмов оператор фермионной ветви.
\end{nogocriterion}

\begin{protocol}
\begin{enumerate}
 \item каноническая составная связность: \textbf{построена};
 \item новый коэффициент: \textbf{нет};
 \item дальний градиент ежа: \textbf{компенсирован};
 \item гладкий профиль $Q,A_Q$: \textbf{построен};
 \item три энергетических интеграла: \textbf{конечны};
 \item полный локальный $SO(3)$-закон: \textbf{нет};
 \item независимые калибровочные бозоны: \textbf{не выведены};
 \item хопфова стабилизаторная линия на границе: \textbf{есть};
 \item гладкое продолжение спинорного накрытия через ядро: \textbf{не доказано};
 \item абсолютные масса и радиус: \textbf{не получены};
 \item следующая проверка:
 \textbf{Фредгольмовость составного оператора Каллиаса};
 \item рождение материи: \textbf{не закрыто}.
\end{enumerate}
\end{protocol}

Машинный сертификат сохранён в
\path{s2t_v6_gauged_projective_spin_cover_parent_gate_results.json}.